\documentclass[11pt]{article}
\usepackage[margin=1.1in]{geometry}
\usepackage{amsmath,amssymb,amsthm}
\usepackage[colorlinks=true,linkcolor=blue,citecolor=blue]{hyperref}

\newtheorem{theorem}{Theorem}[section]
\newtheorem{lemma}[theorem]{Lemma}
\newtheorem{corollary}[theorem]{Corollary}
\newtheorem{remark}[theorem]{Remark}

\newcommand{\R}{\mathbb{R}}
\newcommand{\E}{\mathbb{E}}
\newcommand{\PP}{\mathbb{P}}
\newcommand{\emp}{\bar\mu_N}
\newcommand{\vol}{\operatorname{vol}}

\title{The edge exponent of the uniform doubling condition, II:\\
the matching upper bound}
\author{Yitzchak Shmalo\\
\small Einstein Institute of Mathematics, The Hebrew University of Jerusalem, Jerusalem, Israel\\
\small \texttt{yitzchak.shmalo@gmail.com}}
\date{June 2026}

\begin{document}
\maketitle

\begin{abstract}
In the companion note [\emph{The edge exponent of the uniform doubling
condition and an interpolation conjecture}] we proved that for i.i.d.\ uniform
samples on the unit ball truncated by a single transversal half-space, the
uniform doubling condition of [Shmalo, \emph{Ann.\ Math.\ Artif.\ Intell.}
(2024)] fails with probability $1-o(1)$ at every scaling factor below
$c(\log N/N)^{2/(n+2)}$, and we conjectured the interpolation law:
critical exponent $2/(2n-j)$ when the lowest transversal boundary stratum of
the truncation regions has dimension $j$.  Here we prove the matching upper
bound: the condition holds with probability $1-1/N$ at scaling
$\ell_N=O\bigl((\log N/N)^{2/(n+2)}\bigr)$.  Hence the critical scaling for
the half-space truncation of the ball is exactly
$\Theta\bigl((\log N/N)^{2/(n+2)}\bigr)$, and the interpolation law is now
proved at all three anchor points $j=n-1$, $j=n-2$ and $j=0$.  The geometric
heart is a uniform section bound for $G=B^n\cap\{x_1\le0\}$: at depth $h$
below the extreme point of $G$ in \emph{any} direction, the level section has
$(n-1)$-volume at least $c_nh^{n/2}$.  Its proof splits into one half-section
estimate for feasible directions and, for infeasible directions, two
overlapping regimes about the angular scale $\sin\varphi\asymp\sqrt h$,
each elementary.
\end{abstract}

\section{Setting and result}\label{sec:intro}

Let $K=B^n$ be the closed unit ball, $n\ge2$, $\mu$ the uniform probability
measure on $K$, $T=\{X_1,\dots,X_N\}$ i.i.d.\ samples with empirical measure
$\emp$, and
\[
H:=\{x:\ x_1\le0\},\qquad G:=K\cap H .
\]
Truncated slabs $S_w=\{x:|\langle x-s,u\rangle|\le\frac w2\}\cap H$ and the
uniform doubling condition (UDC) with constants
$(\kappa,\delta,\sigma,\beta,\ell)$ for the single truncation set
$\{(e_1,0)\}$ are as in \cite{ShmaloUDC,companion2,companion6}.  In
\cite{companion6} we proved the lower bound: with probability $1-o(1)$
(for $n\ge3$) the UDC requires $\ell\ge c(\log N/N)^{2/(n+2)}$.  Here we
complete the picture.

\begin{theorem}\label{thm:main}
Let $n\ge2$.  With the constants
$\kappa=1+150$, $\sigma=1$, $\delta=\delta(n)>0$, $\beta=2$ of the companion
papers (here $\widetilde R=1$ and the thickness radius is
$r_0=\tfrac12$, since $B(-\tfrac12e_1,\tfrac12)\subseteq G$), there are
$C(n)$ and $N_0(n,\kappa)$ such that for all $N\ge N_0$, with probability at
least $1-\frac1N$, the sample satisfies the UDC for the truncation set
$\{(e_1,0)\}$ with scaling factor
\[
\ell_N\;=\;C(n)\,\Bigl(\frac{\log N}{N}\Bigr)^{2/(n+2)} .
\]
Consequently, by \cite{companion6}, the critical scaling for this truncation
is exactly $\Theta\bigl((\log N/N)^{2/(n+2)}\bigr)$, and the interpolation
exponent $\frac{2}{2n-j}$ of \cite[Conjecture 1.2]{companion6} is established
at $j=n-2$, alongside the proven cases $j=n-1$
\cite{companion3,companion4} and $j=0$ \cite{companion3,companion5}.
\end{theorem}

By the proof skeleton of \cite{companion2}, as isolated in
\cite[Section 2]{companion3}, Theorem~\ref{thm:main} reduces to a
deterministic estimate: the doubling inequality holds at all scales
$w\in[\ell_N,\beta)$ on the VC event provided every truncated slab of width
$w$ containing a sample point has widened population mass
$\mu(S_{\kappa w})\ge32u_N$, where $u_N\asymp\log N/N$ is the VC rate; and
this follows from a uniform slab-volume bound.  We therefore prove:

\begin{theorem}[Slab volume]\label{thm:slab}
There are $c(n)>0$ and $L_0(n)>0$ such that for every $x\in G$, every unit
vector $u$, and every $0<L\le L_0$,
\[
\vol\Bigl(G\cap\bigl\{y:\ |\langle y-x,u\rangle|\le L\bigr\}\Bigr)
\;\ge\;c(n)\,L^{\frac{n+2}2} .
\]
\end{theorem}

Granting Theorem~\ref{thm:slab}, the proof of Theorem~\ref{thm:main} is
verbatim the assembly of \cite[Proof of Theorem 1.2]{companion3} (equivalently
\cite[Section 3]{companion4}): the population dichotomy of
\cite{companion2} applies since $G$ contains a $\tfrac12$-ball; the VC
relative deviations apply to the class of truncated slabs; a sample point
$x\in S_w$ gives, via Theorem~\ref{thm:slab} applied with
$L=\min\bigl(\tfrac{(\kappa-1)w}2,L_0\bigr)$ and the inclusion
$\{|\langle y-x,u\rangle|\le L\}\subseteq S'_{\kappa w}(u,s)$,
\[
\mu(S_{\kappa w})\;\ge\;\frac{c(n)}{\vol(B^n)}\,
\min\Bigl(\tfrac{(\kappa-1)w}2,\,L_0\Bigr)^{\frac{n+2}2}
\;\ge\;32\,u_N
\qquad(w\ge\ell_N),
\]
by the definition of $\ell_N$ and, for large $w$, the choice of $N_0$.  The
rest of this paper proves Theorem~\ref{thm:slab}.

\section{The section lemma}\label{sec:section}

For a unit vector $u$ write
$\tau(x):=\langle x,u\rangle$, $\tau_{\max}(u):=\max_G\tau$, and
\[
A_u(h)\;:=\;\vol_{n-1}\Bigl(G\cap\{\tau(x)=\tau_{\max}(u)-h\}\Bigr).
\]
Decompose $u=u_1e_1+u'$ with $u'\perp e_1$, and set
$\sin\varphi:=u_1$ when $u_1>0$.

\begin{lemma}[Section lemma]\label{lem:section}
There are $a(n)>0$ and $h_0(n)\in(0,\tfrac14]$ such that for every unit
vector $u$ and every $0<h\le h_0$,
\[
A_u(h)\;\ge\;a(n)\,h^{\,n/2}.
\]
\end{lemma}

The three exponent regimes behind the proof: for $u_1\le0$ the bound holds at
the better rate $h^{(n-1)/2}$; for $u_1>0$ the extreme sits on the edge (or
face) and the rate $h^{n/2}$ is sharp.

\begin{proof}
\emph{Case I: $u_1\le0$.}  Here the unconstrained maximizer of
$\tau$ over $K$, namely $x=u$, lies in $H$, so $\tau_{\max}=1$ with extreme
point $u$, and the level set of $K$ at depth $h$ is the $(n-1)$-ball
\[
D_h:=K\cap\{\tau=1-h\},
\qquad\text{center }c_h=(1-h)u,\quad\text{radius }r_h=\sqrt{2h-h^2}\ge\sqrt h .
\]
If $u=-e_1$, then $D_h=\{x_1=1-h\}\cdot(-1)$ --- precisely, the level set is
$\{x\in K: x_1=-(1-h)\}\subseteq H$, so $A_u(h)=\vol_{n-1}(D_h)
=\omega_{n-1}r_h^{n-1}$.  Otherwise $e_1$ is not parallel to $u$ and the
trace of $\partial H$ on the level hyperplane is an affine hyperplane of that
hyperplane; the center satisfies $(c_h)_1=(1-h)u_1\le0$, i.e.\ the center
lies in the kept half, and a half-ball through the center is kept:
\[
A_u(h)\;\ge\;\tfrac12\,\omega_{n-1}\,r_h^{\,n-1}
\;\ge\;\tfrac{\omega_{n-1}}2\,h^{\frac{n-1}2}
\;\ge\;\tfrac{\omega_{n-1}}2\,h^{\frac n2},
\]
using $h\le1$.

\emph{Case II: $u_1=\sin\varphi>0$, $u'\ne0$.}  Write $p:=u'/|u'|$, so
$|u'|=\cos\varphi$ and, for $x\in G$,
\[
\tau(x)=\sin\varphi\,x_1+\cos\varphi\,\langle x,p\rangle
\;\le\;0+\cos\varphi\cdot|x|\;\le\;\cos\varphi,
\]
with equality iff $x_1=0$, $x=p$: the extreme is the edge point $p$, and
$\tau_{\max}=\cos\varphi$.  Let
\[
\Lambda(h)\;:=\;\vol\bigl(\{x\in G:\ \tau(x)\ge\cos\varphi-h\}\bigr)
\;=\;\int_0^hA_u(t)\,dt .
\]
We first show
\begin{equation}\label{eq:capmass}
\Lambda(h)\;\ge\;b(n)\,h^{\frac{n+2}2}
\qquad(0<h\le h_0),
\end{equation}
and then convert \eqref{eq:capmass} into the section bound.

\emph{Regime IIa: $\sin^2\varphi\ge h\cos\varphi$.}  Write
$h':=h/\cos\varphi\le1$ and note the regime means
$h'\le\tan^2\varphi$.  Parametrize $x=(-t,x_2,x'')$ in coordinates where
$p=e_2$: the cap is
$\{t\ge0,\ |x|\le1,\ x_2\ge1-h'+t\tan\varphi\}$.  Restrict to
\[
0\le t\le\frac{h'}{4\tan\varphi},
\qquad
|x''|^2\le\frac{h'}2 .
\]
Then $t\tan\varphi\le\frac{h'}4$ and
$t^2\le\frac{h'^2}{16\tan^2\varphi}\le\frac{h'}{16}$ by the regime
condition, so, using $\sqrt{1-s}\ge1-s$,
\[
\sqrt{1-t^2-|x''|^2}-\bigl(1-h'+t\tan\varphi\bigr)
\;\ge\;h'-t\tan\varphi-t^2-|x''|^2
\;\ge\;h'\Bigl(1-\tfrac14-\tfrac1{16}-\tfrac12\Bigr)
\;=\;\tfrac{3h'}{16},
\]
i.e.\ the $x_2$-fibre has length at least $\frac{3h'}{16}$ on the whole
restricted $(t,x'')$-box.  Hence
\[
\Lambda(h)\;\ge\;\frac{3h'}{16}\cdot\frac{h'}{4\tan\varphi}\cdot
\omega_{n-2}\Bigl(\frac{h'}2\Bigr)^{\frac{n-2}2}
\;=\;b_1(n)\,\frac{h'^{\frac{n+2}2}}{\tan\varphi}
\;=\;b_1(n)\,h^{\frac{n+2}2}\cdot\frac{\cos^{-\frac n2}\varphi}{\sin\varphi}
\;\ge\;b_1(n)\,h^{\frac{n+2}2},
\]
since $\cos^{-n/2}\varphi\ge1$ and $\sin\varphi\le1$.

\emph{Regime IIb: $\sin^2\varphi\le h\cos\varphi$.}  Then
$1-\cos\varphi\le\sin^2\varphi\le h$, so with
$\hat h:=h+(1-\cos\varphi)\le2h$ the cap contains
$\bigl\{x\in K:\ \langle x,p\rangle\ge1-\tfrac{\hat h}{16},\ x_1\le0\bigr\}$:
indeed for such $x$,
$|x-p|^2=2(1-\langle x,p\rangle)\le\tfrac{\hat h}8$ gives
$|x_1|\le\sqrt{\hat h/8}$, whence
\[
\tau(x)\;\ge\;-\sin\varphi\sqrt{\tfrac{\hat h}8}+\cos\varphi
\Bigl(1-\tfrac{\hat h}{16}\Bigr)
\;\ge\;\cos\varphi-\Bigl(\sqrt{h\cos\varphi}\cdot\sqrt{\tfrac{2h}8}
+\tfrac{2h}{16}\Bigr)
\;\ge\;\cos\varphi-h,
\]
using $\sin\varphi\le\sqrt{h\cos\varphi}$, $\hat h\le2h$ and
$\tfrac h2+\tfrac h8\le h$.  The set
$\{\langle x,p\rangle\ge1-\tfrac{\hat h}{16}\}\cap K$ is a spherical cap
centered at $p$, and $p_1=0$, so the reflection $x_1\mapsto-x_1$ maps it to
itself and exactly half of it lies in $\{x_1\le0\}$.  Hence
\[
\Lambda(h)\;\ge\;\tfrac12\,\vol\Bigl(\text{cap of depth }\tfrac{\hat h}{16}\Bigr)
\;\ge\;b_2(n)\,h^{\frac{n+1}2}\;\ge\;b_2(n)\,h^{\frac{n+2}2},
\]
by the standard two-sided cap-volume asymptotic and $h\le1$.  The two regimes
cover all $\varphi\in(0,\tfrac\pi2)$, proving \eqref{eq:capmass}.

\emph{From cap mass to sections.}  By the Brunn--Minkowski inequality,
$A_u^{1/(n-1)}$ is concave on the range of $\tau$ over $G$, hence
$A_u$ is unimodal in the depth variable $h$.  We first record a bulk anchor.
The inner ball $B^-:=B(-\tfrac12e_1,\tfrac12)\subseteq G$ has $\tau$-center
$m^-:=\tau(-\tfrac12e_1)=-\tfrac12\sin\varphi$, and its sections at
$\tau$-distance at most $\tfrac14$ from $m^-$ have radius at least
$\sqrt{\tfrac14-\tfrac1{16}}=\tfrac{\sqrt3}4$.  The corresponding depths are
$h=\tau_{\max}-\lambda$ with $\lambda\in[m^--\tfrac14,m^-+\tfrac14]$, i.e.\
$h\in[h^\dagger-\tfrac14,h^\dagger+\tfrac14]$ where
$h^\dagger:=\tau_{\max}-m^-=\cos\varphi+\tfrac12\sin\varphi
\in[\tfrac12,\tfrac98]$ for $\varphi\in(0,\tfrac\pi2)$.  In particular there
is a depth $h^\dagger\ge\tfrac12\ge h_0$ with
$A_u(h^\dagger)\ge\omega_{n-1}\bigl(\tfrac{\sqrt3}4\bigr)^{n-1}=:a_1(n)$.

Now let $h\le h_0$ and let $h^*$ be a maximizer of $A_u$.  If $h\le h^*$,
then $A_u$ is nondecreasing on $[0,h]$ and
\[
A_u(h)\;\ge\;\frac1h\int_0^hA_u\;=\;\frac{\Lambda(h)}h
\;\ge\;b(n)\,h^{n/2}
\]
by \eqref{eq:capmass}.  If $h>h^*$, then $A_u$ is nonincreasing on
$[h^*,h^\dagger]\ni h$, so
$A_u(h)\ge A_u(h^\dagger)\ge a_1(n)\ge a_1(n)h^{n/2}$.  In both cases
$A_u(h)\ge a(n)h^{n/2}$ with $a(n):=\min\{b(n),a_1(n)\}$.

\emph{Case II-degenerate: $u'=0$, i.e.\ $u=e_1$.}  Then
$\tau_{\max}=0$, attained on the whole face $K\cap\{x_1=0\}$, and
$A_u(h)=\vol_{n-1}(K\cap\{x_1=-h\})=\omega_{n-1}(1-h^2)^{\frac{n-1}2}
\ge a(n)$ for $h\le\tfrac12$.  (The case $u=-e_1$ was covered in Case I.)
\end{proof}

\section{The slab-volume bound}\label{sec:slab}

\begin{proof}[Proof of Theorem \ref{thm:slab}]
Fix $x\in G$, $u$, and $L\le L_0:=\tfrac{h_0}4$.  Let
$\tau_{\max}:=\tau_{\max}(u)$ and $\tau_{\min}:=\min_G\tau=-\tau_{\max}(-u)$,
and let $A:=A_u$ denote the section profile, extended over the full range
$[\tau_{\min},\tau_{\max}]$, which has length at least $1$ (the inner ball
$B(-\tfrac12e_1,\tfrac12)$ has $\tau$-range of length $1$).  By
Lemma~\ref{lem:section} applied to $u$ and to $-u$,
\[
A(\tau_{\max}-a)\;\ge\;a(n)\Bigl(\tfrac L4\Bigr)^{n/2}=:A^*,
\qquad
A(\tau_{\min}+a')\;\ge\;A^*,
\]
with anchor depths $a=a'=\tfrac L4$.  By unimodality (Brunn--Minkowski),
$A\ge A^*$ on the middle interval
$M:=[\tau_{\min}+\tfrac L4,\ \tau_{\max}-\tfrac L4]$, whose complement in the
range has length $\tfrac L2$.  The slab's level interval
$[\tau(x)-L,\tau(x)+L]$ meets the range in an interval of length at least
$L$ (the range has length $\ge1\ge2L$ and contains $\tau(x)$), hence meets
$M$ in a set of length at least $\tfrac L2$.  Therefore
\[
\vol\Bigl(G\cap\{|\langle y-x,u\rangle|\le L\}\Bigr)
=\int_{\tau(x)-L}^{\tau(x)+L}A
\;\ge\;\frac L2\,A^*
\;=\;\frac{a(n)}{2\cdot4^{n/2}}\,L^{\frac{n+2}2}.
\]
\end{proof}

\section{Conclusion}\label{sec:concl}

Theorem~\ref{thm:main} follows, as explained in Section~\ref{sec:intro}, by
inserting Theorem~\ref{thm:slab} into the proof skeleton of
\cite{companion2,companion3}.  Together with
\cite{companion6}, the critical scaling of the uniform doubling condition for
the half-space truncation of the ball is
$\Theta\bigl((\log N/N)^{2/(n+2)}\bigr)$, and the interpolation law
$\ell_N\asymp(\log N/N)^{2/(2n-j)}$ is now proved at its three anchor strata:
smooth ($j=n-1$), transversal edge ($j=n-2$), and conical point ($j=0$).  The
remaining cases $1\le j\le n-3$ --- transversal strata of intermediate
dimension, produced by $n-1-j$ half-spaces in general position --- should
follow from the same two-step scheme: a section lemma with rate
$h^{(2n-j-2)/2}$ via the regime split at the angular scale
$\sqrt h$ in each transverse direction, and the multi-cap second moment of
\cite{companion4,companion5,companion6} for the lower bound.  We leave the
general case to future work.

\subsection*{Acknowledgments}
The research presented in this paper was supported by the European Research
Council (ERC) under the European Union's Horizon Europe research and
innovation programme (grant agreement No.\ 101041711), by the Israel Science
Foundation (grant number 2258/19) and by the Israel Science Foundation (ISF
Grant 4101/25).

\subsection*{Declaration on generative AI}
During the preparation of this work the author used generative AI tools to
accelerate drafting and revision.  All mathematical claims, proofs, and
bibliographic information were subsequently reviewed and edited by the author,
who takes full responsibility for the content of the manuscript.

\begin{thebibliography}{9}

\bibitem{ShmaloUDC}
Y.~Shmalo,
\emph{Stability of accuracy for the training of DNNs via the uniform doubling
condition},
Ann.\ Math.\ Artif.\ Intell.\ (2024), DOI 10.1007/s10472-023-09919-1.

\bibitem{companion2}
Y.~Shmalo,
\emph{The uniform doubling condition holds with high probability for
independent samples}, manuscript, 2026.

\bibitem{companion3}
Y.~Shmalo,
\emph{The critical scaling of the uniform doubling condition: smooth bodies
versus truncations}, manuscript, 2026.

\bibitem{companion4}
Y.~Shmalo,
\emph{The logarithm is necessary: exact critical scaling of the uniform
doubling condition on the ball}, manuscript, 2026.

\bibitem{companion5}
Y.~Shmalo,
\emph{The exact critical scaling of the uniform doubling condition with
truncations}, manuscript, 2026.

\bibitem{companion6}
Y.~Shmalo,
\emph{The edge exponent of the uniform doubling condition and an
interpolation conjecture}, manuscript, 2026.

\end{thebibliography}

\end{document}
